Fake-news detectors that score near-perfectly within one dataset often fail on another, and artifacts such as publisher markers, document length and named entities are commonly blamed. How much of the gap they explain is rarely measured, and a gap that shrinks after a cue is removed does not show that the cue was the cause. We audit three public English news corpora (ISOT, McIntire, WELFake): under our normalized matching every cleaned ISOT and McIntire document has a WELFake match, so we use a non-overlapping WELFake residual as the third corpus. A factorial intervention design (source-marker masking, length matching, entity masking) with size-matched random-deletion controls, balanced samples, group-aware splits and eight seeds is applied to TF-IDF logistic regression, DistilBERT and RoBERTa. For TF-IDF, removing all three cue families lowers the pooled cross-dataset AUC gap from 0.213 to 0.160 (about 25\% of the gap), driven by length matching and entity masking; entity masking also beats its matched random control after length matching (+0.020 AUC). For both transformers, masking raises transfer but deleting the same number of random words does at least as well, so targeted removal does not outperform matched random deletion; random deletion alone raises RoBERTa's cross-dataset AUC by 0.102 but lowers TF-IDF's by 0.012. Effects depend strongly on the dataset pair: with length matching, the transformers beat their random controls only on transfers between ISOT and McIntire. The tested cue families thus remove part of the TF-IDF gap, and how much a cue seems to matter depends on the model and the dataset pair.
